\subsection{Plancherel 定理}

我们用 A(G) 表示 \(L^{1}(G)\) 的由诸函数 \(f * g\)（\(f, g \in L^{1}(G) \cap L^{2}(G)\)）生成的向量子空间。

命题 7.　空间 A(G) 是 L¹(G) 的一个自伴理想。它含于 L¹(G) ∩ L²(G)，也含于 C(G) 在 L¹(G) 中的像内。

设 \(f \in \mathrm{L}^1(\mathrm{G})\)。对每个 \(g \in \mathrm{L}^2(\mathrm{G})\)，有 \(f * g \in \mathrm{L}^1(\mathrm{G})\)。因此空间 \(\mathrm{L}^1(\mathrm{G}) \cap \mathrm{L}^2(\mathrm{G})\) 是 \(\mathrm{L}^1(\mathrm{G})\) 的一个理想，空间 \(\mathrm{A}(\mathrm{G})\) 亦然。理想 \(\mathrm{A}(\mathrm{G})\) 是自伴的。

设 \(f\) 与 \(g\) 属于 \(\mathrm{L}^2(\mathrm{G})\)。于是卷积 \(f * g\) 是由下式给出的连续函数的类

\[
y \mapsto \int_ {\mathrm{G}} f (y x ^ {- 1}) g (x) d x
\]

第二个断言由此得出。

由于 \(\chi(f*g) = (\chi f)*( \chi g)\) 对 \(\chi \in \widehat{\mathrm{G}}\)、\(f \in \mathrm{L}^1(\mathrm{G})\) 与 \(g \in \mathrm{L}^1(\mathrm{G})\) 成立（INT，VIII，§3，第 1 节，命题 6），我们有 \(\chi h \in \mathrm{A}(\mathrm{G})\) 对所有 \(h \in \mathrm{A}(\mathrm{G})\) 与 \(\chi \in \widehat{\mathrm{G}}\) 成立。由于 \(\varepsilon_x * f = \gamma(x)f\)，且线性表示 \(\gamma\) 在所有 \(\mathrm{L}^p(\mathrm{G})\)（对每个 \(p\)）上都是等距的（INT，VIII，§2，第 5 节，命题 8），我们有 \(\varepsilon_x * f \in \mathrm{A}(\mathrm{G})\) 对所有 \(x \in \mathrm{G}\) 与 \(f \in \mathrm{A}(\mathrm{G})\) 成立。

我们用 \(\widehat{\mathrm{A}}(\mathrm{G})\) 表示 \(\mathrm{A}(\mathrm{G})\) 在 Fourier 变换下的像。它是 \(\mathcal{C}_0(\widehat{\mathrm{G}})\) 的一个子空间。

命题 8.　存在一个滤子基 \(\mathfrak{B}\)（在 \(\mathrm{A}(\mathrm{G}) \cap \mathcal{K}_{+}(\mathrm{G})\) 上），使得下列条件满足：

\begin{enumerate}
\def\labelenumi{(\roman{enumi})}
\item
  对每个这样的 \(\varphi\)（\(\mathfrak{B}\) 的一个成员的元素），有 \(\|\varphi\|_1 = 1\) 且 \(\|\mathcal{F}(\varphi)\|_{\infty} \leqslant 1\)；
\item
  我们有
\end{enumerate}

\[
\lim _ {\varphi , \mathfrak {B}} \varphi \cdot d x = \varepsilon_ {e}
\]

在 G 上具有紧支集的测度所成的空间 \(\mathcal{C}'(\mathrm{G})\) 中成立，该空间赋有在 \(\mathcal{C}(\mathrm{G})\) 的紧子集上一致收敛的拓扑；

\begin{enumerate}
\def\labelenumi{(\roman{enumi})}
\setcounter{enumi}{2}
\tightlist
\item
  我们有
\end{enumerate}

\[
\lim _ {\varphi , \mathfrak {B}} \mathcal {F} (\varphi) = 1
\]

对 \(\widehat{\mathbf{G}}\) 上的紧收敛拓扑成立；

\begin{enumerate}
\def\labelenumi{(\roman{enumi})}
\setcounter{enumi}{3}
\tightlist
\item
  对 \(p = 1\) 或 \(p = 2\)，以及对每个 \(f \in \mathrm{L}^p(\mathrm{G})\)，有 \(\varphi * f \in \mathrm{A}(\mathrm{G})\) 对每个 \(\varphi\)（属于 \(\mathfrak{B}\) 的某个元素）成立，且
\end{enumerate}

\[
\lim _ {\varphi , \mathfrak {B}} \varphi * f = f
\]

在 L \(^{p}\) (G) 中成立。

设 \(K_{0}\) 为 G 中 e 的一个固定的紧邻域。设 \(B_{0}\) 为 G 中 e 的邻域滤子的一个基，由含于 \(K_{0}\) 的紧对称邻域组成（参见 TG，III，第 4 页）。对 \(K \in B_{0}\)，设 \(X_{K}^{\prime}\) 为诸函数 \(\psi \in \mathcal{K}_{+}(G)\)（满足 \(\operatorname{Supp}(\psi) \subset K\) 且 \(\int \psi(x)dx = 1\)）所成的集合；它非空（II，第 200 页，引理 1）。设 \(X_{K}\) 为诸函数 \(\psi * \psi\)（\(\psi \in X_{K}^{\prime}\)）所成的集合。它非空且含于 \(A(G) \cap \mathcal{K}_{+}(G)\)。集合 \(\mathfrak{B}\)（其元素为诸集合 \(X_{K}\)，其中 K 取遍 \(B_{0}\)）是 \(A(G) \cap \mathcal{K}_{+}(G)\) 上的一个滤子基。我们证明 \(\mathfrak{B}\) 满足所需的性质。

若 \(X \in \mathfrak{B}\) 且 \(\varphi \in X\)，则有 \(\| \varphi \|_1 = \int_{G} \varphi(x) dx = 1\)，从而 \(\| \mathcal{F}(\varphi) \|_\infty \leqslant 1\)，这就建立了性质 (i)。

性质 (ii) 由 INT，VIII，§2，第 7 节，引理 4 的推论 1 得出。\(\widehat{\mathbf{G}}\) 的紧子集是 \(\mathcal{C}(\mathbf{G})\) 的紧子集，因而 (ii) 蕴含 \(\lim_{\varphi, \mathfrak{B}} \mathcal{F}(\varphi) = 1\)（对 \(\widehat{\mathbf{G}}\) 上的紧收敛拓扑），即 (iii)。

最后，设 \(p=1\) 或 \(p=2\)。设 \(f\in\mathrm{L}^{p}(\mathrm{G})\)。我们有 \(\varphi*f\to f\) 在 \(\mathrm{L}^{p}(\mathrm{G})\) 中沿滤子 \(\mathfrak{B}\) 成立（INT，VIII，§4，第 7 节，命题 20）。此外，对 \(B_{0}\) 中每个 K 与 \(\varphi\in X_{K}\)，存在 \(\psi\in X_{K}^{\prime}\) 使 \(\varphi=\psi*\psi\)，从而 \(\varphi*f=\psi*(\psi*f)\)。我们有 \(\psi\in L^{1}(G)\cap L^{2}(G)\) 且 \(\psi*f\in L^{1}(G)\cap L^{2}(G)\)，从而 \(\varphi*f\in A(G)\)。

推论 1.　空间 A(G) 在 L \(^{1}\) (G) 与 L \(^{2}\) (G) 中稠密。它在 Stell(G) 中也稠密，并且它在 Fourier 变换下的像 \(\widehat{\mathrm{A}}(G)\) 在 \(\mathcal{C}_{0}(\widehat{\mathrm{G}})\) 中稠密。

命题的断言 (iv) 给出第一个断言。由于 L \(^{1}\) (G) 在 Stell(G) 中稠密，第二个断言由此得出，而最后一个断言进而由 II，第 209 页，命题 5 得出。

推论 2.　对 \(f \in \mathrm{A}(\mathrm{G})\)，设 \(\Omega_f\) 为由 \(\chi \in \widehat{\mathrm{G}}\)（满足 \(\mathcal{F}(f)(\chi) \neq 0\)）所成的集合。诸集合 \(\Omega_f\) 构成 \(\widehat{\mathrm{G}}\) 的一个开覆盖。

这由前一个推论得出，因为对每个 \(\chi \in \widehat{\mathrm{G}}\)，映射 \(f \mapsto \mathcal{F}(f)(\chi)\) 是 \(\mathrm{L}^1(\mathrm{G})\) 的一个非零特征标。

回忆 Stell(G) 的左正则表示 \(\pmb{\gamma}\)（在 \(\mathrm{L}^2 (\mathrm{G})\) 上；参见 I，第 125 页，第 13 节）记作 \(\pmb{\gamma}(\varphi)f = \varphi *f\)，其中 \(\varphi \in \mathrm{Stell}(\mathrm{G})\)，\(f\in \mathrm{L}^2 (\mathrm{G})\)。

引理 3.　对每个 \(f \in \mathrm{A}(\mathrm{G})\)，存在唯一的有界测度 \(\mu_{f}\)（在 \(\widehat{\mathrm{G}}\) 上），使得

\[
(\varphi * f) (e) = \int_ {\widehat {G}} \mathcal {F} (\varphi) d \mu_ {f}\tag{17}
\]

对所有 \(\varphi\in\mathrm{Stell}(\mathrm{G})\) 成立。

此外，对所有 \(f\) 与 \(g\)（在 A(G) 中），我们有等式

\[
\mathcal {F} (f) \cdot \mu_ {g} = \mathcal {F} (g) \cdot \mu_ {f},\tag{18}
\]

它是以 \(\mathcal{F}(f)\) 为密度（关于 \(\mu_g\)）的测度与以 \(\mathcal{F}(g)\) 为密度（关于 \(\mu_f\)）的测度之间的等式。

设 \(f\) 与 \(g\) 为 \(\mathrm{L}^{1}(\mathrm{G}) \cap \mathrm{L}^{2}(\mathrm{G})\) 的元素。对每个 \(\varphi \in \mathrm{Stell}(\mathrm{G})\)，有 \(\varphi * f \in \mathrm{L}^{2}(\mathrm{G})\) 且 \(\| \varphi * f \|_{2} \leqslant \| \varphi \|_{*} \| f \|_{2}\)（I，第 126 页，公式 (8)）。此外，我们有 \(\varphi * (f * g) = (\varphi * f) * g\)（同前，公式 (9)）。后一函数属于 \(\mathcal{C}_{0}(\mathrm{G})\)，即 \(\mathcal{K}(\mathrm{G})\) 在 \(\mathcal{C}(\mathrm{G})\) 中的闭包（INT，VIII，§4，第 5 节，命题 15）。此外，我们有

\[
\left\| \varphi * (f * g) \right\| _ {\infty} \leqslant \left\| \varphi * f \right\| _ {2} \| g \| _ {2} \leqslant \left\| \varphi \right\| _ {*} \| f \| _ {2} \| g \| _ {2}.
\]

由于诸函数 \(f * g\)（\(f\) 与 \(g\) 属于 \(\mathrm{L}^1(\mathrm{G}) \cap \mathrm{L}^2(\mathrm{G})\)）生成 \(\mathrm{A}(\mathrm{G})\)，由此推出 \(\varphi * f \in \mathcal{C}_0(\mathrm{G})\) 对所有 \(f \in \mathrm{A}(\mathrm{G})\) 与 \(\varphi \in \mathrm{Stell}(\mathrm{G})\) 成立，且映射 \(\varphi \mapsto (\varphi * f)(e)\) 是 Stell(G) 上的一个连续线性形式。由于 \(\mathcal{F}\) 是从 Stell(G) 到 \(\mathcal{C}_{0}(\widehat{\mathrm{G}})\) 上的同构（II，第 209 页，命题 5），第一个断言由此得出。

现设 \(f\) 与 \(g\) 属于 \(\mathrm{A}(\mathrm{G})\)。对 \(\varphi \in \mathrm{L}^1 (\mathrm{G})\) 有

\[
\begin{array}{r l r} & & {\big (\mathcal {F} (f) \cdot \mu_ {g} \big) (\mathcal {F} (\varphi)) = \int_ {\widehat {\mathrm{G}}} \mathcal {F} (\varphi) \mathcal {F} (f) d \mu_ {g} = \int_ {\widehat {\mathrm{G}}} \mathcal {F} (\varphi * f) d \mu_ {g}} \\ & & {= ((\varphi * f) * g) (e).} \end{array}\tag{19}
\]

由于 \((\varphi*f)*g=(\varphi*g)*f\)，且 \(L^{1}(G)\) 在 Fourier 变换下的像在 \(\mathcal{C}_{0}(\widehat{\mathrm{G}})\) 中稠密（II，第 210 页，推论），由公式 (19) 推出公式 (18) 对所有 \(f\) 与 \(g\)（在 \(A(G)\) 中）成立。

引理 4.　存在唯一的测度 \(\nu\)（在 \(\widehat{\mathbf{G}}\) 上），使得

\[
\mu_ {f} = \mathcal {F} (f) \cdot \nu
\]

对所有 \(f \in \mathrm{A}(\mathrm{G})\) 成立。对 \(f \in \mathrm{A}(\mathrm{G})\)，有 \(\mathcal{F}(f) \in \mathrm{L}^{1}(\widehat{\mathrm{G}}, \nu) \cap \mathrm{L}^{2}(\widehat{\mathrm{G}}, \nu)\)。设 \(f \in \mathrm{A}(\mathrm{G})\)。用 \(\Omega_{f}\) 表示 \(\widehat{G}\) 中由 \(\chi \in \widehat{G}\)（满足 \(\mathcal{F}(f)(\chi) \neq 0\)）组成的开集。设 \(\varphi\) 为 \(\widehat{G} - \Omega_{f}\) 的特征函数。于是对每个 \(g \in \mathrm{A}(\mathrm{G})\) 有

\[
\int_ {\widehat {\mathrm{G}}} \mathcal {F} (g) d (\boldsymbol {\varphi} \cdot \boldsymbol {\mu} _ {f}) = \int_ {\widehat {\mathrm{G}}} \varphi \mathcal {F} (f) d \mu_ {g} = 0
\]

这由公式 (18) 得到。

由 II，第 212 页，推论 1，像 \(\widehat{\mathrm{A}}(\mathrm{G})\)（\(\mathrm{A}(\mathrm{G})\) 在 Fourier 变换下的）在 \(\mathcal{C}_0(\widehat{\mathrm{G}})\) 中稠密。于是由前面的公式推出 \(\varphi \cdot \mu_f = 0\)，从而 \(\mu_f\) 集中在 \(\Omega_f\) 上（INT，IV，§4，第 7 节，定义 4）。设 \(\nu_f\) 为 \(\Omega_f\) 上以 \(\mathcal{F}(f)^{-1}\) 为密度关于 \(\mu_f|\Omega_f\) 的测度。

诸集合 \(\Omega_f\)（\(f \in \mathrm{A}(\mathrm{G})\)）构成 \(\widehat{\mathrm{G}}\) 的一个开覆盖（推论 2）。对所有 \(f\) 与 \(g\)（在 \(\mathrm{A}(\mathrm{G})\) 中），公式 (18) 表明 \(\nu_f|(\Omega_f \cap \Omega_g) = \nu_g|(\Omega_f \cap \Omega_g)\)。因此，存在唯一的测度 \(\nu\)（在 \(\widehat{\mathrm{G}}\) 上），使得 \(\nu_f = \nu|\Omega_f\) 对所有 \(f \in \mathrm{A}(\mathrm{G})\) 成立（INT，III，§2，第 1 节，命题 1）。

若 \(f \in \mathrm{A}(\mathrm{G})\)，测度 \(\mu_f\) 与 \(\mathcal{F}(f) \cdot \nu\) 都集中在 \(\Omega_f\) 上，并且它们在 \(\Omega_f\) 上的限制都等于 \(\mathcal{F}(f) \cdot \nu_f\)；因此这两个测度相等。

由于 \(\mu_f\) 是有界测度，Fourier 变换 \(\mathcal{F}(f)\) 属于空间 \(\mathrm{L}^1 (\widehat{\mathrm{G}},\nu)\)。此外，由于 \(\mathcal{F}(f)\) 属于 \(\mathcal{C}_0(\widehat{\mathrm{G}})\)，我们也有 \(\mathcal{F}(f)\in \mathrm{L}^2 (\widehat{\mathrm{G}},\nu)\)。

于是公式 (17) 可写成，对 \(\varphi\in\mathrm{Stell}(\mathrm{G})\) 与 \(f\in\mathrm{A}(\mathrm{G})\)：

\[
(\varphi * f) (e) = \int_ {\widehat {G}} \mathcal {F} (\varphi) \mathcal {F} (f) d \nu .\tag{20}
\]

特别地，对诸 \(f\) 与 \(g\)（在 \(\mathrm{A}(\mathrm{G})\) 中），有

\[
\int_ {\widehat {\mathrm{G}}} \mathcal {F} (f) \mathcal {F} (g) d \nu = (f * g) (e) = \int_ {\mathrm{G}} f (x) g (x ^ {- 1}) d x.\tag{21}
\]

命题 9.　引理 4 所刻画的测度 \(\nu\) 是 \(\widehat{\mathbf{G}}\) 上的一个 Haar 测度。

设 \(\chi \in \widehat{\mathrm{G}}\)。对诸 \(f\) 与 \(g\)（在 \(\mathrm{A}(\mathrm{G})\) 中），把公式 (21) 应用于 \(\chi f\) 与 \(\chi g\)。右边不变，故

\[
\nu (\mathcal {F} (f) \mathcal {F} (g)) = \nu (\mathcal {F} (\chi f) \mathcal {F} (\chi g)).
\]

由 II，第 208 页，公式 (12) 与 II，第 208 页，公式 (13)，可得

\[
\nu (\mathcal {F} (f) \mathcal {F} (g)) = \nu (\varepsilon_ {\chi} * (\mathcal {F} (f) \mathcal {F} (g))).
\]

然后由 INT，VIII，§4，第 3 节，命题 7 推出

\[
\nu (\mathcal {F} (f) \mathcal {F} (g)) = (\varepsilon_ {\chi^ {- 1}} * \nu) (\mathcal {F} (f) \mathcal {F} (g)),
\]

即

\[
(\mathscr {F} (f) \cdot \nu) (\mathscr {F} (g)) = (\mathscr {F} (f) \cdot (\varepsilon_ {\chi^ {- 1}} * \nu)) (\mathscr {F} (g)).
\]

由于空间 \(\widehat{\mathrm{A}}(\mathrm{G})\) 在 \(\mathcal{C}_{0}(\widehat{\mathrm{G}})\) 中稠密（推论 1），由此得等式

\[
\mathscr {F} (f) \cdot \nu = \mathscr {F} (f) \cdot (\varepsilon_ {\chi^ {- 1}} * \nu).
\]

因此测度 \(\nu\) 与 \(\varepsilon_{\chi^{-1}} * \nu\) 在开集 \(\Omega_f\)（\(\mathcal{F}(f)\) 在其上非零）上一致。由推论 2 以及 INT，III，§2，第 1 节，命题 1 的推论，这两个测度相等。

这表明测度 \(\nu\) 与 \(\widehat{\mathrm{G}}\) 上的某个 Haar 测度成比例。设 \(f \in \mathrm{A}(\mathrm{G})\)。在公式 (21) 中取 \(g = \widetilde{f}\)。它于是写成

\[
\int_ {\widehat {\mathrm{G}}} | \mathcal {F} (f) | ^ {2} d \nu = \int_ {\mathrm{G}} | f | ^ {2} d x,\tag{22}
\]

这就证明测度 \(\nu\) 非零。因此测度 \(\nu\) 是 \(\widehat{G}\) 上的一个 Haar 测度。

定义 4.　命题 9 中给出的 Haar 测度 \(\nu\)（在 \(\widehat{G}\) 上）称为 G 上给定的 Haar 测度 dx 的对偶测度。

我们常用 \(d\chi\) 或 \(d\widehat{x}\) 表示 \(\widehat{G}\) 上与 Haar 测度 dx 对偶的 Haar 测度。

注.　设 \(a\) 为实数 \(>0\)。若把 \(dx\) 换成测度 \(a \cdot dx\)，则函数 \(f\) 与 \(g \in \mathrm{L}^{1}(\mathrm{G})\) 的卷积换成 \(a(f*g)\)。我们已看到（II，第 209 页，注），\(\mathcal{F}(f)\) 换成 \(a\mathcal{F}(f)\)。因此 \(\mu_{f}\) 不变，而 \(\nu\) 换成 \(a^{-1} \cdot \nu\)。特别地，测度 \(dx \otimes d\hat{x}\)（在 \(\mathrm{G} \times \hat{\mathrm{G}}\) 上）与 \(\mathrm{G}\) 上 Haar 测度的选取无关。

引理 5.　空间 A(\(\hat{G}\)) 在 L²(\(\hat{G}\)) 中稠密。

设 \(h\) 为 \(\mathrm{L}^2(\widehat{\mathrm{G}})\) 中与 \(\widehat{\mathrm{A}}(\mathrm{G})\) 正交的一个元素。对诸 \(f\) 与 \(g\)（在 \(\mathrm{A}(\mathrm{G})\) 中），有 \(\mathcal{F}(f) \cdot \mathcal{F}(g) = \mathcal{F}(f*g) \in \widehat{\mathrm{A}}(\mathrm{G})\)，从而 \(h \cdot \overline{\mathcal{F}(f)}\) 与 \(\mathcal{F}(g)\) 正交。于是，对每个 \(f \in \mathrm{A}(\mathrm{G})\)，函数 \(h \cdot \mathcal{F}(f)\) 与 \(\widehat{\mathrm{A}}(\mathrm{G})\) 正交。但 \(h \cdot \mathcal{F}(f) \in \mathrm{L}^1(\widehat{\mathrm{G}})\)，且 \(\widehat{\mathrm{A}}(\mathrm{G})\) 在 \(\mathcal{C}_0(\widehat{\mathrm{G}})\) 中稠密，故测度 \(h\mathcal{F}(f) \cdot \nu\) 为零，即 \(h\mathcal{F}(f)\) 是 \(\nu\)-局部可忽略的（INT，V，§5，第 3 节，命题 3 的推论 2）。特别地，\(h\) 是 \(\nu\)-局部可忽略的，在集合 \(\Omega_f\)（即诸特征标 \(\chi\)（\(\mathcal{F}(f)(\chi) \neq 0\)）的集合）上。由推论 2，\(h\) 是 \(\nu\)-局部可忽略的，从而为零，因为 \(h\) 属于 \(\mathrm{L}^2(\widehat{\mathrm{G}})\)。证明完毕。

定理 1（Plancherel）.　Fourier 变换在 L²(G) 的子空间 A(G) 上的限制唯一地扩张为从 L²(G) 到 L²(Ĝ) 上的一个等距 Φ。

此外，若 \(f \in \mathrm{L}^{1}(\mathrm{G}) \cap \mathrm{L}^{2}(\mathrm{G})\)，则其 Fourier 变换属于 \(\mathrm{L}^{2}(\widehat{\mathrm{G}})\)，且在 \(\mathrm{L}^{2}(\widehat{\mathrm{G}})\) 中与 \(\Phi(f)\) 一致。

由公式 (22)，\(\mathcal{F}\) 在 A(G) 上的限制是子空间 A(G)（L \(^2\) (G) 的）到子空间 \(\hat{\mathrm{A}}(\mathrm{G})\)（L \(^2\) (\(\hat{\mathrm{G}}\)) 的）上的一个等距。由于 A(G) 在 L \(^2\) (G) 中稠密（II，第 212 页，推论 1），Fourier 变换唯一地扩张为等距 \(\Phi\)，它把 L \(^2\) (G) 映到 L \(^2\) (\(\hat{\mathrm{G}}\)) 的某个闭子空间上。但由于其像包含 \(\hat{\mathrm{A}}(\mathrm{G})\)（它在 L \(^2\) (\(\hat{\mathrm{G}}\)) 中稠密；引理 5），映射 \(\Phi\) 是满的。

现设 \(f \in \mathrm{L}^{1}(\mathrm{G}) \cap \mathrm{L}^{2}(\mathrm{G})\)；我们来证明其 Fourier 变换属于 \(\mathrm{L}^{2}(\widehat{\mathrm{G}})\)。由 II，第 211 页，命题 8, (iv) 以及 \(\mathrm{A}(\mathrm{G})\) 是 \(\mathrm{L}^{1}(\mathrm{G})\) 的理想这一事实，存在 \(\mathrm{A}(\mathrm{G})\) 上的一个滤子基 B，它在 \(\mathrm{L}^{1}(\mathrm{G})\) 与 \(\mathrm{L}^{2}(\mathrm{G})\) 中都收敛于 f。于是我们有

\[
\Phi (f) = \lim _ {g, \mathfrak {B}} \Phi (g) = \lim _ {g, \mathfrak {B}} \mathcal {F} (g)
\]

在 \(\mathrm{L}^{2}(\widehat{\mathrm{G}})\) 中成立，且 \(\mathcal{F}(f)=\lim_{g,\mathfrak{B}}\mathcal{F}(g)\) 在 \(\mathcal{C}_{0}(\widehat{\mathrm{G}})\) 中成立。因此存在一个序列 \((g_{n})\)（在 \(\mathrm{A}(\mathrm{G})\) 中），使 \(\mathcal{F}(g_{n})\) 收敛于 \(\Phi(f)\) 在 \(\mathrm{L}^{2}(\widehat{\mathrm{G}})\) 中，且收敛于 \(\mathcal{F}(f)\) 在 \(\mathcal{C}_{0}(\widehat{\mathrm{G}})\) 中。由 INT，IV，§3，第 4 节，定理 3 及其推论 1，有 \(\mathcal{F}(f)=\Phi(f)\)，特别地 \(\mathcal{F}(f)\in\mathrm{L}^{2}(\widehat{\mathrm{G}})\)。定理证毕。

我们仍用 \(\mathcal{F}\) 表示定理 1 中定义的从 \(\mathrm{L}^2 (\mathrm{G})\) 到 \(\mathrm{L}^2 (\widehat{\mathrm{G}})\) 上的等距，并称之为 \(\mathrm{L}^2 (\mathrm{G})\) 上的 Fourier 变换。同样，Fourier 余变换也唯一地等距扩张到 \(\mathrm{L}^2 (\mathrm{G})\) 上，仍称为 Fourier 余变换，记作 \(\overline{\mathcal{F}}\)。

推论.　假设 G 是紧的且 dx 是 G 上的规范化 Haar 测度。则 G 的酉特征标族是 L²(G) 的一组标准正交基。

由于 G 是紧的，G 的特征标属于 \(L^{2}(G)\)。对诸 \(\chi\) 与 \(\xi\)（在 \(\widehat{G}\) 中），有

\[
\int_ {\mathrm{G}} \overline {{\langle \chi , x \rangle}} \langle \xi , x \rangle d x = \mathcal {F} _ {\mathrm{G}} (d x) (\chi \xi^ {- 1}),
\]

因此 G 的特征标族是正交规范的（II，第 210 页，命题 6）。它还是完全的，因为 \(\chi \in \widehat{\mathbf{G}}\) 与 \(f \in \mathrm{L}^2(\mathrm{G})\) 的内积等于 \(\mathcal{F}_{\mathrm{G}}(f)(\chi)\)，从而 \(f\) 与 \(\widehat{\mathbf{G}}\) 正交当且仅当 \(\mathcal{F}_{\mathrm{G}}(f)\) 在 \(\mathrm{L}^2(\widehat{\mathbf{G}})\) 中为零，当且仅当 \(f\) 为零（定理 1）。

注.　1) 有关 \(\mathrm{L}^{1}(\mathrm{G})\) 上 Fourier 变换的一些公式可以推广到 \(\mathrm{L}^{2}(\mathrm{G})\) 上的 Fourier 变换。特别地，对 \(f \in \mathrm{L}^{2}(\mathrm{G})\) 与 \(\chi \in \widehat{\mathbf{G}}\)，有

\[
\overline {{\mathcal {F}}} (f) = (\chi \mapsto \mathcal {F} (f) (\chi^ {- 1})) = \overline {{\mathcal {F} (\overline {{f}})}}
\]

\[
\mathcal {F} (\varepsilon_ {x} * f) = \eta (x ^ {- 1}) \mathcal {F} (f), \qquad \overline {{\mathcal {F}}} (\varepsilon_ {x} * f) = \eta (x) \mathcal {F} (f)
\]

\[
\mathcal {F} (\chi f) = \varepsilon_ {\chi} * \mathcal {F} (f), \quad \overline {{\mathcal {F}}} (\chi f) = \varepsilon_ {\chi} * \mathcal {F} (f).
\]

若 \(\sigma\) 是 G 的自同构，\(\Delta\) 是 \(\sigma\) 的模（INT，VII，§1，第 4 节，定义 4），则对 \(f \in \mathrm{L}^2(\mathrm{G})\) 有

\[
\mathcal {F} (f \circ \sigma) = \Delta^ {- 1} \mathcal {F} (f) \circ \widehat {\sigma} ^ {- 1}
\]

在 L \(^{2}\) (G) 中成立。

\begin{enumerate}
\def\labelenumi{\arabic{enumi})}
\setcounter{enumi}{1}
\tightlist
\item
  公式
\end{enumerate}

\[
\| \overline {{\mathcal {F}}} (f) \| ^ {2} = \| f \| ^ {2}\tag{23}
\]

（\(f\in \mathrm{L}^2 (\mathrm{G})\)），或

\[
\int_ {\mathrm{G}} f (x) \overline {{g (x)}} d x = \int_ {\widehat {\mathrm{G}}} \mathcal {F} (f) (\chi) \overline {{\mathcal {F} (g) (\chi)}} d \chi\tag{24}
\]

（\(f\) 与 \(g\) 属于 \(\mathrm{L}^2 (\mathrm{G})\)）称为"Plancherel 公式"。

命题 10.　设 \(n \geqslant 0\) 为整数，\(\mathrm{G}_{1}, \cdots, \mathrm{G}_{n}\) 为局部紧交换群。设 \(\mu_{j}\)（\(1 \leqslant j \leqslant n\)）为 \(\mathrm{G}_{j}\) 上的 Haar 测度。设 G 为诸群 \(\mathrm{G}_{j}\)（\(1 \leqslant j \leqslant n\)）的乘积群。设 \(\beta\) 为 II，第 206 页，命题 2 中给出的从 \(\widehat{\mathrm{G}}\) 到 \(\prod \widehat{\mathrm{G}}_{j}\) 上的同构。\(\widehat{\mathrm{G}}\) 上与 G 上的乘积 Haar 测度 \(\mu = \mu_{1} \otimes \cdots \otimes \mu_{n}\) 对偶的 Haar 测度等同于诸 Haar 测度 \(\widehat{\mu}_{j}\) 的乘积。

事实上，对每个族 \((f_{j})\)（由 \(\mathcal{L}^{2}(\mathrm{G}_{j})\) 的非零元素组成），G 上由 \((x_{j})\mapsto\Pi f_{j}(x_{j})\) 定义的函数 f 属于 \(\mathcal{L}^{2}(\mathrm{G})\)，且满足

\[
\int_ {\mathrm{G}} | f | ^ {2} d \mu = \prod_ {j} \int_ {\mathrm{G} _ {j}} | f _ {j} | ^ {2} d \mu_ {j} = \prod_ {j} \int_ {\widehat {\mathrm{G}} _ {j}} | \mathcal {F} _ {\mathrm{G} _ {j}} (f _ {j}) | ^ {2} d \widehat {\mu} _ {j},
\]

这由 Plancherel 公式得出，它证明诸 \(\widehat{\mu}_{j}\) 的乘积 Haar 测度等同于 \(\mu\) 的对偶 Haar 测度。

